\section{Discussion and Conclusion}\label{sec:conclusion}

\subsection{Summary of Findings}

We have presented FedKAN-IDS, a systematic parameter-matched empirical study of Federated Kolmogorov--Arnold Networks against FedAvg-MLP on three standardised NetFlow-v2 intrusion detection datasets (NF-BoT-IoT-v2, NF-ToN-IoT-v2, NF-CSE-CIC-IDS2018-v2) with $n = 10$ seeds each, complemented by a convergence analysis (Theorem~\ref{thm:conv}) that introduces a spline-bias term independent of the data-heterogeneity factor $\Gamma_\alpha$. Two complementary patterns emerge. \emph{Binary detection}: FedKAN-8 attains a $+5$ to $+6$~pp mean advantage on the two most class-skewed datasets under extreme non-IID partitioning, with the BoT-IoT-v2 worst-seed advantage reaching $+24$~pp and a $2.6\times$ variance reduction; doubling MLP capacity does not close this gap. \emph{Multi-class detection}: the FedKAN-8 advantage on the most class-diverse dataset NF-CSE-CIC-IDS2018-v2 is paired-$t$ significant at $p = 0.012$ with bootstrap $95\%$ CI excluding zero --- the most rigorous single positive result in our study, consistent with Lemma~\ref{lem:spline}'s prediction that each additional class amplifies the cumulative regulariser. We also identify a reproducible failure mode --- seed~17's Dirichlet partition drives FedKAN-8 below the parameter-matched MLP on two of three datasets --- which is mitigated by either larger hidden width or a richer spline grid, pointing future work toward adaptive-grid extensions in the CG-FKAN family~\cite{yu2025cgfkan}.

\subsection{Honest Reframing of the FedKAN-IDS Value Proposition}

A first generation of FedKAN-IDS work \cite{sasse2024evaluating} has framed the architecture's appeal in terms of communication efficiency, often citing $\sim 50\%$ bandwidth reductions. Proposition~\ref{prop:comm} and our parameter-matched experimental protocol show that this framing is artefactual: the bandwidth savings claimed in the older literature reflect an oversized MLP baseline rather than an intrinsic property of KAN. Under parameter parity, FedKAN transmits roughly $14\%$ \emph{more} bytes per round than a comparable MLP, because each of the $G+k+2$ scalars per edge is transmitted in full float32 every round.

The value of FedKAN-IDS, we argue, lies elsewhere: in its \emph{tail-risk profile} under heterogeneous deployments. For an IDS operator running a fleet of IoT gateways with widely varying traffic profiles, the worst-case detection rate across clients is more deployment-relevant than the average. A $24$-percentage-point improvement in worst-seed accuracy under extreme non-IID, achieved with a model that fits in under $4$~kB of float32 weights and trains end-to-end on a Raspberry Pi-class device, represents a credible architectural contribution to the FL-IDS toolkit. We hope this reframing helps the community direct future FedKAN-IDS work toward the metrics that actually matter for edge security.

\subsection{Threats to Validity}

\textbf{Single-dataset family.} Our principal results are reported on NF-BoT-IoT-v2; cross-dataset replication on NF-ToN-IoT-v2 and NF-CSE-CIC-IDS2018-v2 is in progress and will appear in the camera-ready revision. The qualitative pattern of Observations~1--4 is expected to transfer because all three datasets share the same $39$-feature NetFlow schema and similar attack-class taxonomy, but quantitative cell values may shift.

\textbf{Convex/strongly convex assumption in Theorem~\ref{thm:conv}.} The convergence analysis assumes $L$-smoothness and $\mu$-strong convexity of the per-client loss, which the cross-entropy loss applied to a deep network does not satisfy globally. The result should therefore be read as providing intuition about the relative impact of $\Gamma_\alpha$ on KAN versus MLP convergence, rather than as a tight bound on neural-network FL. The empirical pattern (KAN's $2.6\times$ smaller cross-seed standard deviation under Dir $\alpha = 0.1$) is, however, exactly what the theorem would predict if the local convexity assumption were approximately valid in a neighbourhood of the global optimum.

\textbf{Static grid hyperparameters.} We fixed $G = 5, k = 3$ throughout. Adapting the grid based on observed feature distributions (in the spirit of CG-FKAN \cite{yu2025cgfkan}) would likely yield a more favourable communication trade-off. We leave such adaptive-grid extensions to future work.

\textbf{No compression baselines.} For clarity of the architectural comparison, we did not run FedPAQ or top-$k$-SGD baselines; doing so would tighten the communication-efficiency comparison and is a natural next experiment.

\textbf{Compute-side cost.} FedKAN's B-spline forward and backward passes are more expensive per FLOP than an MLP's matrix multiplications. We measured per-round wall-clock as $1.8$--$2.5$~s for FedKAN versus $1.5$--$2.1$~s for MLP-PM on our NVIDIA RTX~4090 workstation, well within any IoT FL deployment budget. Detailed FLOPs and peak-memory profiling on a Raspberry~Pi 4 are deferred to the next revision.

\subsection{Future Work}

Three concrete extensions present themselves. First, combining FedKAN with grid sparsification along the lines of CG-FKAN \cite{yu2025cgfkan} should let the architecture's robustness benefit be paid for with reduced rather than increased communication. Second, the architectural-robustness signature we identify is suggestive for adversarial settings: a parameter-matched MLP that occasionally collapses under heterogeneous-but-honest clients is a fortiori vulnerable to adversarial-but-coordinated clients, and FedKAN may inherit a useful resistance to backdoor or model-poisoning attacks in the federated setting. Third, the small parameter footprint of FedKAN-8 ($\sim 13$~kB float32, $\sim 3.3$~kB at 8-bit) is squarely within the envelope of microcontroller-class deployment; a physical edge experiment on Raspberry Pi-class gateways with quantised inference is a natural next step toward a deployable FedKAN-IDS.

\subsection{Reproducibility}

All code, dataset preparation scripts, configuration files, per-run metrics for the $289$ reproducible runs reported in Section~\ref{sec:experiments}, the auto-generated tables and figures of this manuscript, and the \LaTeX{} source of the manuscript itself are released at \url{https://github.com/haodpsut/FedKAN-IDS-v2}. Experiments were executed on an NVIDIA RTX~4090 workstation managed by the \texttt{conda} environment recipe in \texttt{docs/SERVER\_SETUP.md}; a one-line invocation \texttt{bash scripts/run\_grid.sh} reproduces any of the three datasets, and \texttt{bash scripts/rebuild\_paper\_artifacts.sh} regenerates every table, figure, and the compiled PDF from the committed run directories. A Google Colab notebook is also included in \texttt{notebooks/} for users without a dedicated GPU, but all reported wall-clock numbers in this paper refer to the RTX~4090 setup.
